% GENERATED FILE. Edit canonical lessons, then run npm run generate:books.
% canonical-source-hash: fnv1a64:fa8aa310d9672a09
% canonical-lessons: TA-C553-tarkalika, TA-C553-nirantara, TA-C553-inriyamaiyata, TA-C553-marruvali, TA-C553-torayamaka

\chapter{For Now or For Good}
\label{ch:ta-for-now-or-for-good}
\hlchaptermodality{553}

\begin{chapteropening}
\textbf{By the end of this chapter:} \emph{I can say whether something is temporary, permanent or essential.}

Complete the last lesson of chapter 553: I can say whether something is temporary, permanent or essential.
\end{chapteropening}

\section[\texorpdfstring{\ta{தற்காலிக} (taṟkālika)}{தற்காலிக (taṟkālika)}]{\ta{தற்காலிக} (taṟkālika) — temporary}

\label{lesson:TA-C553-tarkalika}



\begin{quote}
Before the new one: say the Tamil for to speed up, then the Tamil for time to act.
\end{quote}

\subsection*{You'll want to know: \ta{தற்காலிக}}
\textbf{\ta{தற்காலிக}} — \emph{taṟkālika} — \textquotedblleft{}temporary\textquotedblright{}.

\begin{quote}
\ta{இது} \ta{ஒரு} \ta{தற்காலிக} \ta{வேலை}, \ta{மூன்று} \ta{மாதங்கள்} \ta{மட்டுமே}.
\end{quote}

\begin{quote}
\emph{itu oru taṟkālika vēlai, mūṉṟu mātaṅkaḷ maṭṭumē.}
\end{quote}

\begin{quote}
This is a temporary job, only for three months.
\end{quote}

It stands before a noun. As an adverb: \ta{தற்காலிகமாக}, for the time being.

\begin{cousinweb}
\textbf{\ta{தற்காலிக}} (\emph{taṟkālika}) is borrowed from a Sanskrit adjective built on \emph{tatkāla}, \textquotedblleft{}happening at that same time, immediately\textquotedblright{} (Monier-Williams).
\end{cousinweb}

\begin{grammarlens}[title={What you've built}]
One more word for saying what happens if.
\end{grammarlens}

\subsection*{Your turn}
\begin{itemize}
\raggedright
  \item \emph{Say it:} \emph{taṟkālika}
  \item \emph{Say it:} \emph{taṟkālika}, once more
  \item \emph{Say it:} \emph{itu oru taṟkālika vēlai, mūṉṟu mātaṅkaḷ maṭṭumē.} — This is a temporary job, only for three months.
  \item \emph{Recall:} say the Tamil for to speed up, then the Tamil for time to act, then say \emph{taṟkālika} again
\end{itemize}

\begin{tcolorbox}[breakable,colback=teal!4,colframe=teal!35!black,title={Before you move on}]
What does \ta{தற்காலிக} mean? (\textquotedblleft{}Temporary\textquotedblright{}.) Say it once more.
\end{tcolorbox}
\section[\texorpdfstring{\ta{நிரந்தர} (nirantara)}{நிரந்தர (nirantara)}]{\ta{நிரந்தர} (nirantara) — permanent}

\label{lesson:TA-C553-nirantara}



\begin{quote}
Before the new one: say the Tamil for time to act, then the Tamil for temporary.
\end{quote}

\subsection*{You'll want to know: \ta{நிரந்தர}}
\textbf{\ta{நிரந்தர}} — \emph{nirantara} — \textquotedblleft{}permanent\textquotedblright{}.

\begin{quote}
\ta{அவருக்கு} \ta{அரசு} \ta{அலுவலகத்தில்} \ta{நிரந்தர} \ta{வேலை} \ta{கிடைத்தது}.
\end{quote}

\begin{quote}
\emph{avarukku aracu aluvalakattil nirantara vēlai kiṭaittatu.}
\end{quote}

\begin{quote}
He got a permanent job in a government office.
\end{quote}

It stands before a noun, like \ta{தற்காலிக} (\emph{taṟkālika}). As an adverb: \ta{நிரந்தரமாக}.

\begin{cousinweb}
\textbf{\ta{நிரந்தர}} (\emph{nirantara}) is borrowed from Sanskrit \emph{nirantara}, \textquotedblleft{}having no interval, uninterrupted, constant\textquotedblright{} (Monier-Williams).
\end{cousinweb}

\begin{grammarlens}[title={What you've built}]
One more word for saying what happens if.
\end{grammarlens}

\subsection*{Your turn}
\begin{itemize}
\raggedright
  \item \emph{Say it:} \emph{nirantara}
  \item \emph{Say it:} \emph{nirantara}, once more
  \item \emph{Say it:} \emph{avarukku aracu aluvalakattil nirantara vēlai kiṭaittatu.} — He got a permanent job in a government office.
  \item \emph{Recall:} say the Tamil for time to act, then the Tamil for temporary, then say \emph{nirantara} again
\end{itemize}

\begin{tcolorbox}[breakable,colback=teal!4,colframe=teal!35!black,title={Before you move on}]
What does \ta{நிரந்தர} mean? (\textquotedblleft{}Permanent\textquotedblright{}.) Say it once more.
\end{tcolorbox}
\section[\texorpdfstring{\ta{இன்றியமையாத} (iṉṟiyamaiyāta)}{இன்றியமையாத (iṉṟiyamaiyāta)}]{\ta{இன்றியமையாத} (iṉṟiyamaiyāta) — essential; indispensable}

\label{lesson:TA-C553-inriyamaiyata}



\begin{quote}
Before the new one: say the Tamil for temporary, then the Tamil for permanent.
\end{quote}

\subsection*{You'll want to know: \ta{இன்றியமையாத}}
\textbf{\ta{இன்றியமையாத}} — \emph{iṉṟiyamaiyāta} — \textquotedblleft{}essential; indispensable\textquotedblright{}.

\begin{quote}
\ta{உணவும்} \ta{தண்ணீரும்} \ta{வாழ்க்கைக்கு} \ta{இன்றியமையாதவை}.
\end{quote}

\begin{quote}
\emph{uṇavum taṇṇīrum vāḻkkaikku iṉṟiyamaiyātavai.}
\end{quote}

\begin{quote}
Food and water are essential for life.
\end{quote}

Word for word \textquotedblleft{}that cannot be done without\textquotedblright{}: \ta{இன்றி} \textquotedblleft{}without\textquotedblright{} + \ta{அமையாத} \textquotedblleft{}not fitting\textquotedblright{}.

\begin{cousinweb}
\textbf{\ta{இன்றியமையாத}} (\emph{iṉṟiyamaiyāta}) is a native phrase grown into one word: \ta{இன்றி} (\emph{iṉṟi}), \textquotedblleft{}without\textquotedblright{}, and \ta{அமையாத}, the negative adjective of \ta{அமை} (\emph{amai}), \textquotedblleft{}to fit, to come about\textquotedblright{} (DEDR 162).
\end{cousinweb}

\begin{grammarlens}[title={What you've built}]
One more word for saying what happens if.
\end{grammarlens}

\subsection*{Your turn}
\begin{itemize}
\raggedright
  \item \emph{Say it:} \emph{iṉṟiyamaiyāta}
  \item \emph{Say it:} \emph{iṉṟiyamaiyāta}, once more
  \item \emph{Say it:} \emph{uṇavum taṇṇīrum vāḻkkaikku iṉṟiyamaiyātavai.} — Food and water are essential for life.
  \item \emph{Recall:} say the Tamil for temporary, then the Tamil for permanent, then say \emph{iṉṟiyamaiyāta} again
\end{itemize}

\begin{tcolorbox}[breakable,colback=teal!4,colframe=teal!35!black,title={Before you move on}]
What does \ta{இன்றியமையாத} mean? (\textquotedblleft{}Essential; indispensable\textquotedblright{}.) Say it once more.
\end{tcolorbox}
\section[\texorpdfstring{\ta{மாற்றுவழி} (māṟṟuvaḻi)}{மாற்றுவழி (māṟṟuvaḻi)}]{\ta{மாற்றுவழி} (māṟṟuvaḻi) — an alternative; a way round}

\label{lesson:TA-C553-marruvali}



\begin{quote}
Before the new one: say the Tamil for permanent, then the Tamil for essential.
\end{quote}

\subsection*{You'll want to know: \ta{மாற்றுவழி}}
\textbf{\ta{மாற்றுவழி}} — \emph{māṟṟuvaḻi} — \textquotedblleft{}an alternative; a way round\textquotedblright{}.

\begin{quote}
\ta{சாலை} \ta{மூடப்பட்டால்} \ta{மாற்றுவழியில்} \ta{போகலாம்}.
\end{quote}

\begin{quote}
\emph{cālai mūṭappaṭṭāl māṟṟuvaḻiyil pōkalām.}
\end{quote}

\begin{quote}
If the road is closed, we can take another route.
\end{quote}

Word for word \textquotedblleft{}a change-way\textquotedblright{}: \ta{மாற்று} (\emph{māṟṟu}) + \ta{வழி} (\emph{vaḻi}).

\begin{cousinweb}
\textbf{\ta{மாற்றுவழி}} (\emph{māṟṟuvaḻi}) is a native compound: \ta{மாற்று} (\emph{māṟṟu}), \textquotedblleft{}a change, a substitute\textquotedblright{}, from \ta{மாறு} (DEDR 4834), and \ta{வழி} (\emph{vaḻi}), \textquotedblleft{}a way\textquotedblright{}.
\end{cousinweb}

\begin{grammarlens}[title={What you've built}]
One more word for saying what happens if.
\end{grammarlens}

\subsection*{Your turn}
\begin{itemize}
\raggedright
  \item \emph{Say it:} \emph{māṟṟuvaḻi}
  \item \emph{Say it:} \emph{māṟṟuvaḻi}, once more
  \item \emph{Say it:} \emph{cālai mūṭappaṭṭāl māṟṟuvaḻiyil pōkalām.} — If the road is closed, we can take another route.
  \item \emph{Recall:} say the Tamil for permanent, then the Tamil for essential, then say \emph{māṟṟuvaḻi} again
\end{itemize}

\begin{tcolorbox}[breakable,colback=teal!4,colframe=teal!35!black,title={Before you move on}]
What does \ta{மாற்றுவழி} mean? (\textquotedblleft{}An alternative; a way round\textquotedblright{}.) Say it once more.
\end{tcolorbox}
\section[\texorpdfstring{\ta{தோராயமாக} (tōrāyamāka)}{தோராயமாக (tōrāyamāka)}]{\ta{தோராயமாக} (tōrāyamāka) — approximately; roughly}

\label{lesson:TA-C553-torayamaka}



\begin{quote}
Before the new one: say the Tamil for essential, then the Tamil for an alternative.
\end{quote}

\subsection*{You'll want to know: \ta{தோராயமாக}}
\textbf{\ta{தோராயமாக}} — \emph{tōrāyamāka} — \textquotedblleft{}approximately; roughly\textquotedblright{}.

\begin{quote}
\ta{கூட்டத்துக்குத்} \ta{தோராயமாக} \ta{ஐம்பது} \ta{பேர்} \ta{வந்தார்கள்}.
\end{quote}

\begin{quote}
\emph{kūṭṭattukkut tōrāyamāka aimpatu pēr vantārkaḷ.}
\end{quote}

\begin{quote}
Roughly fifty people came to the meeting.
\end{quote}

The -\ta{ஆக} adverb of \ta{தோராயம்}, a rough estimate. A near synonym is \ta{கிட்டத்தட்ட} (\emph{kiṭṭattaṭṭa}).

\begin{cousinweb}
The origin of \textbf{\ta{தோராயமாக}} (\emph{tōrāyamāka}) is not settled: it is the -\ta{ஆக} adverb of \ta{தோராயம்} (\emph{tōrāyam}), \textquotedblleft{}a rough estimate\textquotedblright{}, whose source is unclear.
\end{cousinweb}

\begin{grammarlens}[title={What you've built}]
One more word for saying what happens if.
\end{grammarlens}

\subsection*{Your turn}
\begin{itemize}
\raggedright
  \item \emph{Say it:} \emph{tōrāyamāka}
  \item \emph{Say it:} \emph{tōrāyamāka}, once more
  \item \emph{Say it:} \emph{kūṭṭattukkut tōrāyamāka aimpatu pēr vantārkaḷ.} — Roughly fifty people came to the meeting.
  \item \emph{Recall:} say the Tamil for essential, then the Tamil for an alternative, then say \emph{tōrāyamāka} again
\end{itemize}

\begin{tcolorbox}[breakable,colback=teal!4,colframe=teal!35!black,title={Before you move on}]
What does \ta{தோராயமாக} mean? (\textquotedblleft{}Approximately; roughly\textquotedblright{}.) Say it once more.
\end{tcolorbox}
